\section{Discussion}
\label{sec:discussion}

\paragraph{Benchmarking Kubernetes pods from one set of measurements.}
Because the user code is the same whichever way a pod is submitted, one measurement of
the workload can label both a raw Job and a Flyte task, as long as the accounting for the
pod's edges is right. Getting it right took real pods. Sibling processes in one pod
confirmed the arithmetic, but only per-container cgroups showed where tmpfs pages are
charged, that the first release's init containers could not hold their inputs, and that
the Flyte runtime was less than half the size we had assumed. Anyone deriving Kubernetes
labels this way should check them against the cgroup, not against another RSS
measurement.

\paragraph{What a model learns from Jobs.}
A model trained on Jobs reads the manifest. It changes its memory request when the
volume medium changes, and it does so on first runs where nothing else gives the staging
away. Flyte tasks add nothing to that model's accuracy on Jobs, and a Flyte-only model
does not transfer well even after an explicit adapter. For a team that submits pods both
ways, the practical advice is to train on the pods it actually runs. Flyte history is still
useful for sizing Flyte tasks, and the corrected adapter recovers part of the gap when Job
history is missing.

\paragraph{Familiar and unfamiliar code.}
The gap between the two splits is the main open problem. On familiar code the best
approaches are within 7--8\% of the oracle. On new kinds of code, the best trained model is
39\% above it, and a careful rule matches it. Pooling the trees with a language model helps
on familiar code and hurts on new code, because the pool's weight is tuned on familiar
families. The gate recovers the better of the two on each split, but on this benchmark its
detector has an easy job, since new families usually bring new libraries. A benchmark
that separates unfamiliar code from unfamiliar libraries would test it properly.

\paragraph{Cost, not accuracy.}
Accuracy on the label and cost often disagree. Cost-aware decoding over-asks on purpose,
the Flyte-to-Job adapter was judged a failure until its runtime constant was corrected, and
the ranking on new code flips when retries become expensive. A benchmark for this
decision has to score money, and it has to say what a failure costs.

\section{Limitations}
\label{sec:limitations}

The workloads are synthetic. They use real libraries and real code paths, but the tasks
are generated, the data is random, and runs take seconds to minutes. Production tasks are
longer and messier. The seven unseen families were chosen by hand, and six of them
introduce a library that no training family uses, so the benchmark does not test new code
built only from familiar libraries. The people who wrote the generators also wrote the tree
features and the rules, and the task names give away the family, although removing them
changes regret by less than four points.

The labels carry three known biases. The \code{ru\_maxrss} floor and memory-mapped inputs
make about a third of labels one cell larger than strictly needed, so the oracle is
slightly generous. Process pools are labelled by their largest process rather than their
total, which is why two of 100 real pods failed at their label. And RightSizeMem has no
per-second CPU series, so its throttling is modelled from mean cores. The Job labels are
derived rather than measured for every run. We validated them on 100 runs in real pods
and 176 in a single pod, but we did not re-measure GPU runs as Jobs, and the four stagings
leave out sidecars, CSI drivers and other init work.

The cost model uses list prices, with no spot discounts or bin-packing, and prices a retry
at one oracle cost by default; \cref{sec:res-sens} shows how much that choice matters. We
trained three times, which shows that rankings on new code within about ten points are not
stable, but three seeds are not enough to put tight intervals on that variation. The
zero-shot model was scored on a sample of about 300 runs per split, and only once.

\section{Conclusion}
\label{sec:conclusion}

RightSize Bench scores the request a scheduler makes before a Kubernetes pod runs, in
dollars, across 33 families of task code submitted both as raw Jobs and through Flyte.
Checking its labels in real pods fixed two errors and measured the cost of Flyte's runtime
at about 70\,MiB. A small model trained on Jobs reads the manifest and, pooled with
gradient-boosted trees, comes within 7\% of the oracle on familiar code. On unfamiliar
code, the best trained model and a careful rule are both about 39\% above the oracle, and
the rule is safer when failures are expensive. Closing that gap, and testing a novelty
detector on code that looks familiar but is not, is the work that remains.
